\documentclass[12pt, oneside]{article}

\usepackage{xparse}

\usepackage{geometry}
\geometry{letterpaper}
\usepackage[parfill]{parskip}
\usepackage{float}

% Math packages
\usepackage{amssymb}
\usepackage{amsmath}
\usepackage{amsthm}
\usepackage{mathalpha}

% Tables
\usepackage{booktabs}
\usepackage{afterpage}

% Hyperlinks
\usepackage{hyperref}

% References
\numberwithin{equation}{section}
\usepackage[numbers]{natbib}

% TikZ (optional for later)
\usepackage{tikz}
\usetikzlibrary{shapes,arrows,chains}

% Theorem environments
\theoremstyle{definition}
\newtheorem{definition}{Definition}[section]
\newtheorem{proposition}{Proposition}[section]
\newtheorem{axiom}{Axiom}[section]

\theoremstyle{plain}
\newtheorem{theorem}{Theorem}[section]
\newtheorem{lemma}[theorem]{Lemma}
\newtheorem{corollary}[theorem]{Corollary}
\newtheorem{example}[theorem]{Example}

\theoremstyle{remark}
\newtheorem*{remark}{Remark}

\title{Dyadic Brackets and Density-One Descent in the Collatz Map}
\author{Wayne Brassem}

\begin{document}

% Custom commands
\NewDocumentCommand{\setN}{}{\mathbb{N}}
\NewDocumentCommand{\setZ}{}{\mathbb{Z}}
\NewDocumentCommand{\setQ}{}{\mathbb{Q}}

% Document commands which provide hyperlinks to OEIS sequences referenced herein
\NewDocumentCommand{\OEIS}{m}{\href{https://oeis.org/#1}{#1}}

\maketitle

\begin{abstract}

We study the Collatz map via the fully accelerated transformation
\[
f(x)=\frac{3x+1}{2^{v_2(3x+1)}}.
\]

Finite trajectories are encoded by division words, which record the
2-adic valuations encountered along an orbit. Each such word induces
an affine map whose admissibility imposes a congruence constraint on
initial values. This yields a symbolic–affine framework in which
admissible words correspond to residue classes modulo powers of two.

We show that these constraints refine exponentially: a division word
of length $m$ restricts initial values to a residue class modulo a
power of $2$ that is comparable to, but strictly larger than, $3^m$,
while the number of admissible words grows like $\rho^m$ with $\rho < 3$.

This creates a competition between combinatorial growth and arithmetic
refinement: while the number of admissible words grows like $\rho^m$
with $\rho < 3$, each word imposes a congruence constraint modulo a
power of $2$ that grows strictly faster than $3^m$. This imbalance
forces exponential sparsification of admissible initial conditions.

As a consequence, the set of integers that fail to admit a finite
descent prefix has natural density zero. In particular, almost all
integers eventually decrease under iteration of the Collatz map.

\end{abstract}

\newpage
\tableofcontents

% ============================================================
\newpage
\section{Introduction}

The $3x+1$ problem, also known as the Collatz conjecture~\cite{Collatz1937}, concerns
the behavior of the map
\[
T(x) =
\begin{cases}
\displaystyle
    \phantom{3x} \frac{x}{2} & \text {if } x \text{ is even}, \\[8pt]
\displaystyle
    3x+1,                    & \text{if $x$ is odd}.
\end{cases}
\]
on the positive integers. It asserts that every positive integer eventually reaches $1$
under iteration of $T$. Despite its elementary formulation, the conjecture has resisted
proof for decades~\cite{Lagarias2010,Terras1976}. Extensive computational verification
confirms convergence for very large ranges~\cite{OliveiraESilva2010}, but does not
reveal the global structure governing trajectory behavior.

A major advance by Tao~\cite{Tao2019} showed that almost all integers, in the sense of
logarithmic density, eventually decrease under iteration of the Collatz map. This result
is obtained via probabilistic averaging rather than an explicit structural decomposition
of the integers.

The present work develops a deterministic, arithmetic framework for analyzing typical
Collatz behavior based on a symbolic–affine encoding of trajectories. Using this approach,
we encode the dynamics symbolically rather than studying individual trajectories with
the fully accelerated map
\[
f(x)=\frac{3x+1}{2^{v_2(3x+1)}},
\]
which maps odd integers to odd integers by compressing each odd step into a single
transformation.

Each trajectory under $f$ is represented by a finite \emph{division word}, recording the
number of factors of $2$ removed (the 2-adic valuation) at each iterate.

These words admit an exact affine representation obtained by symbolic composition of the
accelerated map: each division word induces a map of the form
\[
F_\omega(x) = \frac{3^m x + c(\omega)}{2^{D(\omega)}},
\]
and determines a corresponding congruence class modulo $2^{D(\omega)}$. Thus division words
determine canonical arithmetic \emph{constraint classes} (residue classes) in the integers,
once realizability is imposed.

This correspondence allows the dynamics to be studied at the level of residue classes
rather than individual orbits. Convergent division words—those for which the associated
affine map is contractive—identify families of integers that decrease after a fixed
number of steps. The global problem is thereby reduced to understanding how these
families are distributed and how their densities accumulate.

Two competing effects govern this distribution. On one hand, the number of admissible
convergent division words grows exponentially with length. On the other hand, each word
imposes a congruence constraint whose modulus grows exponentially with the total 2-adic
valuation. Using structural constraints on admissible division words, encoded by an
increment grammar, we show that the combinatorial growth rate is strictly subcritical:
the number of admissible words grows like $\rho^m$ for some $\rho < 3$, while the moduli
scale like powers of $2$ satisfying $3^m < 2^{D(\omega)} < 2\cdot 3^m$.

This transforms the Collatz problem into a competition between symbolic growth and
arithmetic refinement, with convergence governed by their relative exponential rates.

This imbalance leads to \emph{exponential sparsification}: the proportion of integers
admitting a convergent division word of length $m$ decays exponentially in $m$.
Summing over all lengths then shows that the complement—integers that avoid all such
descent patterns—has natural density zero.

Consequently, the set of integers whose trajectory under iteration of $f$ eventually
decreases has natural density one.

While this does not exclude the existence of exceptional trajectories, any such
exceptions must lie within a set of zero density and cannot arise from any scalable
arithmetic structure generated by the symbolic dynamics.

% \paragraph{Structure of the Paper.}
% The argument proceeds as follows:
% \begin{itemize}
% \item Section~\ref{sec:division-counts} introduces division words and their total division counts.
% \item Section~\ref{sec:affine-representation} develops the affine representation of division words and their congruence structure.
% \item Section~\ref{sec:affine-grammar} establishes admissibility and minimal valuation constraints.
% \item Section~\ref{sec:symbolic-growth} controls the growth of admissible words via the increment grammar.
% \item Section~\ref{sec:affine-constraint-fibers} identifies the residue class structure induced by division words.
% \item Section~\ref{sec:exponential-sparsification} proves exponential sparsification of admissible sets.
% \item Section~\ref{sec:density-one-descent} establishes density-one descent.
% \end{itemize}


% ============================================================
\section{Grammar–Drift–Block Framework}
\label{sec:grammar-drift-block}
% ============================================================

\begin{table}[htbp]
\centering
\caption{Minimum block size, $M$, for logarithmic convexity of $a(n) = A186009(n)$}
\label{tab:A186009}
\vspace{4pt}
\begin{tabular}{@{}rccrr@{}}
\toprule
$M$ & $a_{n} a_{n+2M} \ge a_{n+M}$ & {S/D} & $n$ & $a(n)$ \\
\midrule
1 & $a_{ 1} a_{ 3} \ge a_{ 2} a_{ 2}$ & S &  1 &            1 \\
1 & $a_{ 2} a_{ 4} \ge a_{ 3} a_{ 3}$ & S &  2 &            1 \\
2 & $a_{ 3} a_{ 7} \ge a_{ 5} a_{ 5}$ & S &  3 &            1 \\
1 & $a_{ 4} a_{ 6} \ge a_{ 5} a_{ 5}$ & D &  4 &            2 \\
2 & $a_{ 5} a_{ 9} \ge a_{ 7} a_{ 7}$ & S &  5 &            3 \\
1 & $a_{ 6} a_{ 8} \ge a_{ 7} a_{ 7}$ & D &  6 &            7 \\
1 & $a_{ 7} a_{ 9} \ge a_{ 8} a_{ 8}$ & S &  7 &           12 \\
3 & $a_{ 8} a_{14} \ge a_{11} a_{11}$ & D &  8 &           30 \\
1 & $a_{ 9} a_{11} \ge a_{10} a_{10}$ & D &  9 &           85 \\
2 & $a_{10} a_{14} \ge a_{12} a_{12}$ & S & 10 &          173 \\
1 & $a_{11} a_{13} \ge a_{12} a_{12}$ & D & 11 &          476 \\
1 & $a_{12} a_{14} \ge a_{13} a_{13}$ & S & 12 &          961 \\
4 & $a_{13} a_{21} \ge a_{17} a_{17}$ & D & 13 &         2652 \\
1 & $a_{14} a_{16} \ge a_{15} a_{15}$ & D & 14 &         8045 \\
3 & $a_{15} a_{21} \ge a_{18} a_{18}$ & S & 15 &        17637 \\
1 & $a_{16} a_{18} \ge a_{17} a_{17}$ & D & 16 &        51033 \\
2 & $a_{17} a_{21} \ge a_{19} a_{19}$ & S & 17 &       108950 \\
1 & $a_{18} a_{20} \ge a_{19} a_{19}$ & D & 18 &       312455 \\
1 & $a_{19} a_{21} \ge a_{20} a_{20}$ & S & 19 &       663535 \\
4 & $a_{20} a_{28} \ge a_{24} a_{24}$ & D & 20 &      1900470 \\
1 & $a_{21} a_{23} \ge a_{22} a_{22}$ & D & 21 &      5936673 \\
2 & $a_{22} a_{26} \ge a_{24} a_{24}$ & S & 22 &     13472296 \\
1 & $a_{23} a_{25} \ge a_{24} a_{24}$ & D & 23 &     39993895 \\
1 & $a_{24} a_{26} \ge a_{25} a_{25}$ & S & 24 &     87986917 \\
4 & $a_{25} a_{33} \ge a_{29} a_{29}$ & D & 25 &    257978502 \\
1 & $a_{26} a_{28} \ge a_{27} a_{27}$ & D & 26 &    820236724 \\
3 & $a_{27} a_{33} \ge a_{31} a_{31}$ & S & 27 &   1899474678 \\
1 & $a_{28} a_{30} \ge a_{29} a_{29}$ & D & 28 &   5723030586 \\
2 & $a_{29} a_{33} \ge a_{31} a_{31}$ & S & 29 &  12809477536 \\
1 & $a_{30} a_{32} \ge a_{31} a_{31}$ & D & 30 &  38036848410 \\
1 & $a_{31} a_{33} \ge a_{32} a_{32}$ & S & 31 &  84141805077 \\
4 & $a_{32} a_{40} \ge a_{36} a_{36}$ & D & 32 & 248369601964 \\
1 & $a_{33} a_{35} \ge a_{34} a_{34}$ & D & 33 & 794919136728 \\
\bottomrule
\end{tabular}
\end{table}

This section formalises the interaction between the combinatorial
structure induced by $A020914(n)$ and the induced growth dynamics
of $C(n) = \sum_{k=0}^n N(k)$.

The central idea is that admissible division words generate a
finite-state grammar, which induces a bounded drift automaton.
This in turn yields a uniform block inequality governing residual decay.

% ============================================================
\section{Symbolic Encoding of Collatz Trajectories}
\label{sec:division-counts}

% ------------------------------------------------------------
\subsection{Division Words}

\begin{definition}[Division Word]
A division word of length $m \ge 0$ is a sequence
\[
\omega = (d_0, \dots, d_{m-1}),
\]
where each $d_i \in \mathbb{N}$ is intended to represent the number of factors of $2$
removed at the $i$-th iterate of the accelerated map
\[
f(x) = \frac{3x+1}{2^{v_2(3x+1)}}.
\]
\end{definition}

\begin{remark}
Since $f$ maps odd integers to odd integers, each $3f^i(x)+1$ is even,
so $d_i \ge 1$ for all $i$.
\end{remark}

\begin{definition}[Realization]
An odd integer $x \in \mathbb{N}$ realizes $\omega$ if
\[
d_i = v_2(3f^i(x)+1), \quad 0 \le i < m.
\]
\end{definition}

% ------------------------------------------------------------
\subsection{Total Division Count}

\begin{definition}
The total division count of a word $\omega$ is
\[
D(\omega) := \sum_{i=0}^{m-1} d_i.
\]
\end{definition}

\begin{remark}
The quantity $D(\omega)$ measures the total dyadic contraction
associated with the symbolic path $\omega$.
\end{remark}

% ============================================================
\section{Affine Representation of Division Words}
\label{sec:affine-representation}

\begin{definition}
Each division word $\omega$ of length $m$ induces an associated affine map
\[
F_\omega(x) = \frac{3^m x + c(\omega)}{2^{D(\omega)}},
\]
obtained by symbolic composition of the accelerated map.
\end{definition}

\begin{proposition}
The map $\omega \mapsto F_\omega$ defines a semigroup homomorphism from
concatenation of division words (read left to right) to composition of affine maps.
\end{proposition}

% ------------------------------------------------------------
\subsection{Example}

Consider convergent segments with $m=4$ iterates of $f$.  
The minimal total division count for $m=4$ is
\[
A020914(4) = 7.
\]

One admissible division word attaining this minimum is
\[
\omega = (1,1,1,4),
\]
for which
\[
D(\omega) = 1+1+1+4 = 7.
\]

Composing the accelerated map according to $\omega$ yields
\[
f_\omega(x)=\frac{3^4 x + c(\omega)}{2^7}
           = \frac{81x + 65}{128},
\]
so the canonical affine constant for this word is
\[
c(\omega)=65.
\]

The integrality condition
\[
81x + 65 \equiv 0 \pmod{128}
\]
determines a unique residue class
\[
x \equiv 15 \pmod{128}.
\]
This congruence defines a constraint fiber with modulus $2^7$.

\medskip

\noindent
\textbf{Interpretation.}
The word $\omega = (1,1,1,4)$ means:

\begin{itemize}
    \item At each of the first three iterates of $f$, $3x+1$ contains exactly one factor of $2$.
    \item At the fourth iterate of $f$, $3x+1$ contains four factors of $2$.
\end{itemize}

Under the accelerated map
\[
f(x)=\frac{3x+1}{2^{v_2(3x+1)}},
\]
all of these divisions occur within a single iterate. Thus $d_3=4$ means that
the fourth iterate divides by $2^4$ in that step. The constant $c(\omega)$ is
fixed entirely by this symbolic structure and ensures that precisely the
integers in the residue class $x \equiv 15 \pmod{128}$ produce integer iterates
under the affine form.

This example illustrates that the behavior of trajectories is governed
entirely by the symbolic data encoded in $\omega$ and its associated
affine map. The initial values that realize a given word are precisely
those satisfying the induced congruence constraint, independent of any
particular choice of representative.

\medskip

\noindent
In general, not every division word arises from an actual trajectory;
the condition for realizability will be formulated later as an explicit
congruence constraint, defining the class of \emph{admissible} words.

% ============================================================
\section{Affine Symbolic Admissibility Grammar}
\label{sec:affine-grammar}

\begin{remark}[Division Grammar as a Dynamical Encoding]
Recall from Section~\ref{sec:division-counts} that a \emph{division word}
$\omega=(d_0,\dots,d_{m-1})$ records the exact 2-adic valuation
$d_i = v_2(3 x_i + 1)$, occurring at each iterate of the fully accelerated
Collatz map.

Unlike a purely combinatorial encoding, division words retain essential
dynamical information. Each division word $\omega=(d_0,\dots,d_{m-1})$
determines the sequence of multiplicative factors of the linear coefficient
\[
r_i = \frac{3}{2^{d_i}},
\]
so that the contractive or expansive contribution of each iterate is
preserved in the linear coefficient (with the affine term handled separately
via $c(\omega)$). This structure enables quantitative control of admissible
trajectories via affine growth and to formulate sufficient conditions for
descent, even though individual trajectories are collapsed under the encoding.
\end{remark}

% ------------------------------------------------------------
\subsection{Explicit Form of the Affine Constant}
\begin{lemma}[Explicit Form of the Affine Constant]
Let $\omega = (d_0,\dots,d_{m-1})$ be a division word. Then the affine map
\[
F_\omega(x) = \frac{3^m x + c(\omega)}{2^{D(\omega)}}
\]
has the explicit form
\[
c(\omega)
=
\sum_{j=0}^{m-1}
3^{m-1-j}\cdot 2^{\sum_{i=0}^{j-1} d_i},
\]
where empty sums are defined as $0$.
\end{lemma}

% ------------------------------------------------------------
\subsection{Stepwise Realizability}

\begin{theorem}[Stepwise Realizability of Division Words]
\label{thm:stepwise-realizability}
Let $\omega = (d_0,\dots,d_{m-1})$ be a division word, and define $c(\omega)$ by
\[
c(\omega)
=
\sum_{j=0}^{m-1}
3^{m-1-j}\cdot 2^{\sum_{i=0}^{j-1} d_i}.
\]

Then the identity
\[
2^{D(\omega)} x_m = 3^m x_0 + c(\omega)
\]
holds if and only if the sequence $(x_0,\dots,x_m)$ satisfies the accelerated
Collatz relations
\[
x_{i+1} = \frac{3x_i+1}{2^{d_i}}, \quad 0 \le i < m,
\]
with each intermediate step integral.

In particular, the congruence
\[
3^m x_0 + c(\omega) \equiv 0 \pmod{2^{D(\omega)}}
\]
is equivalent to the existence of a trajectory realizing $\omega$.
\end{theorem}

In particular, the congruence ensures that each intermediate iterate
$x_i$ is an integer, so realizability is enforced at every step.

\begin{remark}[Structural Interpretation]
Each term in $c(\omega)$ corresponds to the additive contribution of a single
``+1'' in the Collatz iteration, propagated forward by subsequent multiplications
by $3$ and backward-scaled by preceding divisions by $2$.
\end{remark}

% ------------------------------------------------------------
\subsection{Admissibility of Division Words}

\begin{definition}[Admissibility]
\label{def:admissible-division-word}
A division word $\omega$ is admissible if the congruence
\[
3^m x + c(\omega) \equiv 0 \pmod{2^{D(\omega)}}
\]
has a solution.
\end{definition}
\begin{remark}
Admissibility is therefore an arithmetic condition on the word $\omega$
itself, expressed as a divisibility constraint on the associated affine
map. When this condition holds, the set of all integers $x$ for which
$F_\omega(x)$ is integral defines a unique congruence constraint modulo
$2^{D(\omega)}$, corresponding to the residue class induced by $\omega$.
\end{remark}

% ------------------------------------------------------------
\subsection{Residue Class Representation of Division Words}

\begin{theorem}[Affine Constraint Representation]
Let $\omega$ be an admissible division word of length $m$.

Then there exists a unique residue $r_\omega \pmod{2^{D(\omega)}}$ such that
\[
x \text{ realizes } \omega \iff x \equiv r_\omega \pmod{2^{D(\omega)}}.
\]

Thus admissible division words correspond to affine congruence constraints
of modulus $2^{D(\omega)}$.
\end{theorem}

\begin{proof}
From the affine representation,
\[
F_\omega(x) = \frac{3^m x + c(\omega)}{2^{D(\omega)}},
\]
integrality is equivalent to
\[
3^m x + c(\omega) \equiv 0 \pmod{2^{D(\omega)}}.
\]

Since $3^m$ is odd, it is invertible modulo $2^{D(\omega)}$, yielding
\[
x \equiv -c(\omega) \cdot (3^m)^{-1} \pmod{2^{D(\omega)}}.
\]

This determines a unique residue modulo $2^{D(\omega)}$.

Conversely, any $x$ in this class satisfies the integrality condition.
By Theorem~\ref{thm:stepwise-realizability}, this is equivalent to the
existence of a trajectory whose successive iterates satisfy
\[
x_{i+1} = \frac{3x_i+1}{2^{d_i}}, \quad 0 \le i < m,
\]
and hence realize the division word $\omega$.

Thus admissible division words are in one-to-one correspondence with affine
congruence constraints modulo powers of two.
\end{proof}

% \begin{lemma}[Prefix Uniqueness of Realizations]
% \label{lem:prefix-uniqueness}
% Let $\omega$ and $\omega'$ be admissible division words.
% If an integer $x$ lies in both cylinders $C_\omega$ and $C_{\omega'}$,
% then one of the words is a prefix of the other.
% \end{lemma}

% \begin{proof}
% An integer $x$ determines a unique sequence of valuations under the
% accelerated map $f$, and hence a unique division word describing its
% trajectory.

% If $x \in C_\omega \cap C_{\omega'}$, then both $\omega$ and $\omega'$
% must agree with the initial segment of this valuation sequence.
% Therefore, one must be a prefix of the other.
% \end{proof}

\begin{lemma}[Prefix Nesting of Cylinders]
\label{lem:prefix-uniqueness}
Let \(\omega\) and \(\omega'\) be admissible division words.

If
\[
C_\omega \cap C_{\omega'} \neq \varnothing,
\]
then one of the words is a prefix of the other.
\end{lemma}

\begin{proof}
Every integer \(x\) determines a unique infinite valuation sequence
\[
(v_2(3x+1),\, v_2(3f(x)+1),\, v_2(3f^2(x)+1),\dots).
\]

If
\[
x \in C_\omega \cap C_{\omega'},
\]
then both \(\omega\) and \(\omega'\) must agree with the initial segment
of this valuation sequence.

Therefore the two words cannot diverge at any position before the shorter
word terminates. Hence one word is a prefix of the other.
\end{proof}

\begin{corollary}[Disjointness at Equal Length]
Distinct admissible division words of the same length define disjoint cylinders.
\end{corollary}

\begin{proof}
If two words of equal length intersected, then by
Lemma~\ref{lem:prefix-uniqueness},
one would be a prefix of the other.
Equal lengths force equality of the words.
\end{proof}

% ------------------------------------------------------------
\subsection{Injectivity of the Affine Encoding}
\begin{lemma}[Injectivity of the Affine Encoding at Fixed Length]
\label{lem:injectivity}
Let $\omega = (d_0,\dots,d_{m-1})$ and $\omega' = (d'_0,\dots,d'_{m-1})$
be division words of the same length $m$. Let
\[
c(\omega)
=
\sum_{j=0}^{m-1}
3^{m-1-j}\cdot 2^{D_j},
\quad
D_j = \sum_{i=0}^{j-1} d_i,
\]
and define $c(\omega')$ analogously.

If $\omega \ne \omega'$, then
\[
c(\omega) \ne c(\omega').
\]

Equivalently, the map $\omega \mapsto c(\omega)$ is injective on admissible
division words of fixed length $m$.
\end{lemma}

\begin{proof}
Suppose $\omega \ne \omega'$, and let
\[
k = \min\{\, j : d_j \ne d'_j \,\}
\]
be the first index at which the two words differ. Then for all $j < k$,
we have $D_j = D'_j$, while $D_k \ne D'_k$.

Consider the difference
\[
\Delta := c(\omega) - c(\omega')
=
\sum_{j=0}^{m-1}
3^{m-1-j}\bigl(2^{D_j} - 2^{D'_j}\bigr).
\]
The terms with $j<k$ cancel, so
\[
\Delta
=
3^{m-1-k}\bigl(2^{D_k} - 2^{D'_k}\bigr)
+
\sum_{j=k+1}^{m-1}
3^{m-1-j}\bigl(2^{D_j} - 2^{D'_j}\bigr).
\]

Let $\delta_k = \min(D_k, D'_k)$. Without loss of generality,
assume $D_k > D'_k$. Then
\[
2^{D_k} - 2^{D'_k}
=
2^{D'_k}(2^{D_k - D'_k} - 1),
\]
where $2^{D_k - D'_k} - 1$ is odd. Hence
\[
v_2(2^{D_k} - 2^{D'_k}) = D'_k = \delta_k.
\]

For $j>k$, both $D_j$ and $D'_j$ strictly exceed $\delta_k$, since the
prefix sums increase by positive integers. Therefore
\[
v_2\!\left(3^{m-1-j}\bigl(2^{D_j} - 2^{D'_j}\bigr)\right)
>
\delta_k
\quad \text{for all } j>k.
\]

It follows that $\Delta$ is a sum of one term with 2-adic valuation
exactly $\delta_k$ and remaining terms of strictly higher 2-adic valuation,
so no cancellation at the lowest 2-adic order is possible. Hence
\[
v_2(\Delta) = \delta_k,
\]
and in particular $\Delta \ne 0$.

Therefore $c(\omega) \ne c(\omega')$, completing the proof.
\end{proof}

This implies distinct division words induce distinct affine constraints.

\begin{corollary}[Distinct Minimal Cylinders at Equal Length]
Let \(\omega\neq\omega'\) be admissible division words of equal length.

Then the corresponding residue classes modulo powers of two are distinct.
\end{corollary}

% ------------------------------------------------------------
\subsection{Minimal Division Count}

\begin{lemma}[Minimal Division Count]
For any admissible division word $\omega=(d_0,\dots,d_{m-1})$ of length $m$,
\[
d_0 + \cdots + d_{m-1} \;\ge\; A020914(m),
\]
where
\[
A020914(m) = \lceil m \log_2(3) \rceil,
\]
the minimal integer $k$ such that $2^k > 3^m$.

Equality holds if and only if $\omega$ achieves the minimal possible total
division count among admissible words of length $m$.
\end{lemma}

\begin{proof}
If
\[
D(\omega) < A020914(m),
\]
then by definition of \(A020914(m)\),
\[
2^{D(\omega)} \le 3^m.
\]

Hence
\[
\frac{3^m}{2^{D(\omega)}} \ge 1,
\]
so the affine coefficient is noncontractive.

Therefore no admissible word can achieve convergence with fewer than
\(A020914(m)\) total divisions.
\end{proof}

% ------------------------------------------------------------
\subsection{Convergent Division Word}

\begin{definition}[Convergent Division Word]
An admissible division word $\omega = (d_0,\dots,d_{m-1})$ is said to be
\emph{convergent} (in the linear sense) if
\[
\frac{3^m}{2^{d_0+\cdots+d_{m-1}}} < 1.
\]
\end{definition}

% ============================================================
\section{Reconstruction of Integers from a Division Word}
\label{sec:reconstruction}

\subsection{Inverse Cylinder Construction}

Let $\omega = (d_0,\dots,d_{m-1})$ be an admissible division word with total division count
\[
D(\omega) = \sum_{i=0}^{m-1} d_i.
\]

From Section~\ref{sec:affine-grammar}, $\omega$ induces the affine map
\[
F_\omega(x) = \frac{3^m x + c(\omega)}{2^{D(\omega)}}.
\]

Admissibility implies the congruence condition
\[
3^m x + c(\omega) \equiv 0 \pmod{2^{D(\omega)}}.
\]

Since $3^m$ is invertible modulo $2^{D(\omega)}$, this yields the unique residue class
\[
x \equiv -c(\omega)\,(3^m)^{-1} \pmod{2^{D(\omega)}}.
\]

\subsection{Canonical Parametrization of the Cylinder}

Let $r_\omega$ denote the canonical representative of this residue class in
$\{0,1,\dots,2^{D(\omega)}-1\}$. Then all integers realizing $\omega$ are given by
\[
x(k) = r_\omega + k\cdot 2^{D(\omega)}, \quad k \in \mathbb{N}_0.
\]

\subsection{Affine Reconstruction Form}

Applying $F_\omega$ to $x(k)$ yields
\[
F_\omega(x(k))
= \frac{3^m (r_\omega + k2^{D(\omega)}) + c(\omega)}{2^{D(\omega)}}.
\]

Expanding,
\[
F_\omega(x(k))
= k\cdot 3^m + \frac{3^m r_\omega + c(\omega)}{2^{D(\omega)}}.
\]

Define the constant
\[
K_\omega := \frac{3^m r_\omega + c(\omega)}{2^{D(\omega)}}.
\]

Then the affine form becomes
\[
F_\omega(x(k)) = k\cdot 3^m + K_\omega.
\]

\subsection{Interpretation}

Thus each division word induces a bijection between:
\begin{itemize}
\item integers in a residue class $x(k)$, and
\item affine images of the form $k\cdot 3^m + K_\omega$.
\end{itemize}

This provides a complete inverse description of the division-word cylinder.

\subsection{Example}
Solving the congruence
\[
3^4 x + 65 \equiv 0 \pmod{128}
\]
gives
\[
81x + 65 \equiv 0 \pmod{128}.
\]

Reducing $81 \equiv -47 \pmod{128}$, we obtain
\[
-47x + 65 \equiv 0 \pmod{128}
\quad \Longrightarrow \quad
47x \equiv 65 \pmod{128}.
\]

Since $47$ is invertible modulo $128$, with
\[
47^{-1} \equiv 79 \pmod{128},
\]
we multiply both sides to obtain
\[
x \equiv 79 \cdot 65 \equiv 15 \pmod{128}.
\]

Thus the division word $\omega = (1,1,1,4)$ generates the residue class
\[
x = 15 + 128k, \quad k \in \mathbb{N}_0.
\]


% ============================================================
\section{Jump-Induced Expansion of the Admissible Grammar}
\label{sec:jump-expansion}
% ============================================================

The admissible division-word grammar is governed entirely by the
prefix-sum constraints induced by
\[
d(n)=A020914(n)=\lfloor n\log_2 3\rfloor+1.
\]

For a division word
\[
\omega=(d_0,\dots,d_{m-1}),
\]
define the prefix sums
\[
H(i)=\sum_{j=0}^{i} d_j.
\]

Admissibility is characterized by the conditions
\[
H(i)<d(i+1)
\quad (0\le i<m-1),
\]
together with the terminal equality
\[
H(m-1)=d(m).
\]

These conditions define a constrained lattice region inside
\[
\mathbb{N}^m,
\]
whose integer points correspond exactly to admissible division words
of length \(m\).

% ------------------------------------------------------------

\subsection{Propagation Versus Expansion}

The increment sequence
\[
\delta(n)=d(n+1)-d(n)
\]
takes values in \(\{1,2\}\).

The distinction between these two increment types governs the
combinatorial evolution of the admissible grammar.

% ------------------------------------------------------------

\begin{lemma}[Unit Increments Preserve the Existing Grammar]
\label{lem:single-no-new}
Suppose
\[
d(n+1)-d(n)=1.
\]

% Then extending the admissible search from length \(n\) to length
% \(n+1\) introduces no new interior admissibility region.
introduces no new interior admissibility directions beyond propagation
of existing prefixes. The admissible words at level \(n+1\) arise solely
from propagation of previously admissible prefix configurations.
\end{lemma}

\begin{proof}[Structural mechanism]
A unit increment increases the terminal admissibility bound by exactly
one:
\[
d(n+1)=d(n)+1.
\]

Since every division word component satisfies \(d_i\ge1\), the final
coordinate is forced to absorb precisely this additional unit.

No new slack is introduced into the earlier prefix constraints
\[
H(i)<d(i+1),
\]
so the admissible lattice region does not acquire any new interior
directions.

Thus the grammar propagates existing admissible configurations but
creates no genuinely new combinatorial branches.
\end{proof}

% ------------------------------------------------------------

\begin{lemma}[Jump Increments Create New Admissible Words]
\label{lem:jump-new-words}
Suppose
\[
d(n+1)-d(n)=2.
\]

Then extending the admissible search from length \(n\) to length
\(n+1\) necessarily creates new admissible division words that did
not occur at earlier lengths.
\end{lemma}

\begin{proof}[Geometric mechanism]
A jump increment enlarges the terminal admissibility bound by two:
\[
d(n+1)=d(n)+2.
\]

After accounting for the mandatory positivity condition
\[
d_i\ge1,
\]
one additional unit of slack remains available.

This additional slack creates new admissible lattice points satisfying
\[
H(i)<d(i+1)
\]
for all intermediate prefixes while still achieving the terminal
constraint
\[
H(n)=d(n+1).
\]

Consequently the admissible lattice region acquires new integer points,
corresponding to genuinely new admissible division words.
\end{proof}

% ------------------------------------------------------------

\subsection{Minimal-Length Uniqueness}

The terminal equality condition
\[
H(m-1)=d(m)
\]
plays a fundamental role in preventing double counting.

% ------------------------------------------------------------

\begin{lemma}[Minimal-Length Uniqueness]
\label{lem:minimal-uniqueness}
Every admissible division word is counted exactly once at its minimal
admissible length.
\end{lemma}

\begin{proof}
The strict inequalities
\[
H(i)<d(i+1)
\quad (0\le i<m-1)
\]
exclude premature termination at shorter lengths.

Meanwhile the terminal equality
\[
H(m-1)=d(m)
\]
forces the total division count to match the minimal admissible dyadic
threshold for length \(m\).

Thus a division word cannot simultaneously satisfy the admissibility
conditions at two different lengths.

Therefore each admissible word appears exactly once in the grammar.
\end{proof}

% ------------------------------------------------------------

\subsection{Enumeration by the Admissible Grammar}

The admissible grammar therefore evolves by two distinct mechanisms:

\begin{enumerate}
\item unit increments propagate existing admissible configurations,
\item jump increments create genuinely new admissible lattice points.
\end{enumerate}

Because the admissibility conditions prevent double counting,
the total number of admissible division words of length \(m\)
is exactly the number of distinct lattice points satisfying the
prefix constraints.

% Consequently,
% \[
% A186009(m+1)
% \]
% counts precisely the admissible division words of length \(m\).

\begin{theorem}[Enumeration Theorem]
\label{thm:enumeration}
The sequence \(A186009(m+1)\) is exactly the number of admissible
division words of length \(m\).
\end{theorem}

Equivalently, \OEIS{A186009} is the exact counting sequence of the
minimal admissible division-word grammar induced by the Beatty
threshold sequence \(A020914\).

\begin{theorem}[First-Passage Characterization]
An admissible division word of length \(m\) is precisely a division word
whose cumulative dyadic exponent satisfies
\[
\frac{3^{i+1}}{2^{H(i)}} > 1
\quad (0\le i<m-1),
\]
and
\[
\frac{3^m}{2^{H(m-1)}} < 1.
\]

Equivalently, admissible words are exactly the minimal first-passage
crossings of the dyadic contraction threshold defined by \(A020914\).
\end{theorem}

% ============================================================
\section{Symbolic Growth and Spectral Bounds}
\label{sec:symbolic-growth}

The combinatorial complexity of admissible division words governs the number of
affine constraint fibers at each depth. This growth is encoded by the sequence
\OEIS{A186009}, which counts convergent division words and hence, by injectivity
of the affine encoding, the associated admissible residue classes by stopping time.

% ------------------------------------------------------------
\subsection{Combinatorial growth model}

The sequence \OEIS{A186009} arises from a horizontal-sum recurrence of Pascal type,
augmented by boundary duplication events governed by the increment grammar
\OEIS{A022921}. This grammar induces a finite symbolic system on two local operations:
single increments ($s$) and double increments ($d$), whose admissible
concatenations generate all configurations counted by \OEIS{A186009}.

The resulting dynamics decompose into a finite set of admissible local patterns,
whose concatenations generate all admissible configurations. The key structural
feature of this grammar is that locally expansive configurations cannot occur
arbitrarily often: they are separated by mandatory contractive patterns.

As a consequence, the total number of admissible configurations exhibits
\emph{submultiplicative growth}, and admits a well-defined exponential growth
rate in the sense of a limsup.

% ------------------------------------------------------------
\section{Denominator Beattic Composition and Jump-Compensated Numerator Growth}

\begin{lemma}[Diminishing Tower Error]
Let $\alpha=\log_2 3$, and let $q_k$ be the ordered sequence of convergent and semiconvergent denominators associated to $\alpha$.

Then
\[
\epsilon_k := q_k\alpha - p_k
\]
for nearest integers $p_k$ satisfies
\[
|\epsilon_k|\to 0.
\]

Hence the corresponding tower blocks become asymptotically exact dyadic/ternary synchronizations.
\end{lemma}


% ------------------------------------------------------------
\section{Beatty Structure and Tower Decomposition of the Denominator}
\label{sec:beatty-towers}

We formalize the macroscopic structure governing the denominator growth
through the Beatty sequence associated with $\log_2 3$.

% ------------------------------------------------------------

\begin{definition}[Dyadic Threshold Sequence]
Let
\[
\alpha := \log_2 3,
\qquad
d(n) := A020914(n) = \lfloor n\alpha \rfloor + 1.
\]
\end{definition}

The increment word
\[
\Delta(n) := d(n+1)-d(n)
\]
takes values in $\{1,2\}$ and defines a Sturmian (Beatty) sequence.

% ------------------------------------------------------------

\begin{definition}[Tower Lengths]
A sequence of integers $\{q_k\}$ is called a tower sequence if
\[
q_k \alpha = p_k + \varepsilon_k,
\]
where $p_k \in \mathbb{Z}$ and $|\varepsilon_k|$ is minimal among all
denominators up to $q_k$.

We refer to $q_k$ as a \emph{tower length} and $p_k$ as its associated
dyadic exponent.
\end{definition}

Empirically, the tower sequence begins as
\[
5,\,7,\,12,\,41,\,53,\,306,\,359,\,665,\,15601,\,16266,\,31867,\,\dots
\]

% ------------------------------------------------------------

\begin{lemma}[Diminishing Approximation Error]
\label{lem:tower-error}
Let $\varepsilon_k := q_k \alpha - p_k$.
Then
\[
|\varepsilon_k| \to 0 \quad \text{as } k \to \infty.
\]
\end{lemma}

\begin{proof}
This follows from the theory of continued fractions and best rational
approximations of irrational numbers.
\end{proof}

% ------------------------------------------------------------

\begin{lemma}[Alternating Approximation Sign]
\label{lem:alternating}
The errors $\varepsilon_k$ alternate in sign.  That is,
tower approximations alternate between underestimates and overestimates
of $\alpha$.
\end{lemma}

\begin{proof}
This is a standard property of convergents of continued fractions.
\end{proof}

% ------------------------------------------------------------

\begin{definition}[Expansive and Contractive Towers]
A tower length $q_k$ is called:
\begin{itemize}
\item \emph{expansive} if $\varepsilon_k < 0$ (i.e.\ $q_k\alpha < p_k$),
\item \emph{contractive} if $\varepsilon_k > 0$.
\end{itemize}
\end{definition}

Thus expansive towers slightly undercount dyadic growth, while
contractive towers slightly overcount it.

% ------------------------------------------------------------

\begin{lemma}[Tower Partition Structure]
\label{lem:tower-partition}
Tower lengths occur in finite clusters of the form
\[
\{5,7,12\},\quad
\{41,53\},\quad
\{306,359,665\},\quad
\{15601,16266,31867\},\ \dots
\]
in which:
\begin{itemize}
\item the smallest element is expansive,
\item all remaining elements are contractive.
\end{itemize}
\end{lemma}

\begin{proof}[Heuristic justification]
These clusters arise from convergents and semiconvergents of $\alpha$.
The smallest element in each cluster lies below $\alpha$ (negative error),
while subsequent refinements overshoot (positive error) before the next
cycle begins.
\end{proof}

% ------------------------------------------------------------

\begin{lemma}[Exact Jump Count on Towers]
\label{lem:exact-jumps}
Let $q_k$ be a tower length.  Then the increment word over the interval
$[n,n+q_k)$ contains an exact number of double steps (jumps), independent
of $n$.
\end{lemma}

\begin{proof}[Sketch]
Since
\[
d(n+q_k)-d(n) \in \{p_k, p_k+1\},
\]
and $\varepsilon_k$ is small, the total number of increments equal to $2$
is fixed across the block.  The Sturmian structure ensures that all
length-$q_k$ factors have identical Parikh counts.
\end{proof}

% ------------------------------------------------------------

\begin{lemma}[Terminal Double--Double Structure]
\label{lem:dd-terminal}
Every tower block ends in a double--double pattern $(2,2)$.
\end{lemma}

\begin{proof}[Heuristic]
The atomic building blocks of the increment word are the $5$-cycle and
$7$-cycle, both of which terminate in $(2,2)$.  Since all tower lengths
are concatenations of these atomic cycles, the terminal structure is
preserved.
\end{proof}

% ------------------------------------------------------------

\begin{theorem}[Macroscopic Denominator Law]
\label{thm:denominator-law}
Let $q_k$ be a tower sequence.  Then over any interval of length $q_k$:
\begin{enumerate}
\item The total dyadic exponent satisfies
\[
d(n+q_k)-d(n) = p_k + O(1),
\]
with $|O(1)| \le 1$.

\item The jump count is exact and independent of $n$.

\item The multiplicative imbalance satisfies
\[
\frac{3^{q_k}}{2^{p_k}} = 2^{\varepsilon_k},
\]
where $\varepsilon_k \to 0$.

\item Expansive towers ($\varepsilon_k<0$) create slack, while
contractive towers ($\varepsilon_k>0$) consume slack.

\item The sequence of towers alternates between expansion and contraction,
with diminishing magnitude.
\end{enumerate}
\end{theorem}

% ------------------------------------------------------------

\begin{remark}[Denominator Rigidity]
The denominator is completely governed by the Beatty structure of
$\log_2 3$.  All macroscopic irregularities reduce to a controlled,
alternating sequence of small multiplicative errors across tower scales.
\end{remark}


% ============================================================
\section{Canonical Minimal Envelope and Regenerative Growth}
\label{sec:minimal-envelope}
% ============================================================

Let
\[
a(n)=A186009(n),
\]
and let
\[
d(n)=A020914(n)=\lfloor n\log_2 3\rfloor+1.
\]

Define the increment sequence
\[
\delta(n):=d(n+1)-d(n)\in\{1,2\}.
\]

A value $\delta(n)=2$ induces a terminal duplication event in the
admissible tuple grammar.

The central structural observation is that the admissible grammar is
monotone with respect to the placement of duplication events:
earlier duplication permanently increases future tuple mass.

Consequently, the canonical Beatty phase beginning at intercept $0$
produces the minimal admissible growth envelope.

% ------------------------------------------------------------

\subsection{Jump Schedules and Tuple Evolution}

% ------------------------------------------------------------

\begin{definition}[Jump Schedule]
A jump schedule is a sequence
\[
J=(\delta_0,\delta_1,\delta_2,\dots),
\qquad
\delta_i\in\{1,2\},
\]
describing the placement of terminal duplication events in the
tuple recursion.

The canonical schedule is the Beatty/Sturmian increment sequence
\[
\delta_i=d(i+1)-d(i),
\]
associated to
\[
d(i)=\lfloor i\log_2 3\rfloor+1.
\]
\end{definition}

% ------------------------------------------------------------

\begin{definition}[Generating Tuple]
Let
\[
T_n=(t_1,t_2,\dots,t_r)
\]
denote the admissible generating tuple at level $n$.

The tuple evolves by:
\begin{enumerate}
\item cumulative Pascal-type addition,
\item terminal duplication whenever $\delta(n)=2$.
\end{enumerate}

Define the total tuple mass by
\[
|T_n|:=\sum_i t_i.
\]
\end{definition}

% ------------------------------------------------------------

\subsection{Monotonicity of the Grammar}

% ------------------------------------------------------------

\begin{lemma}[Positivity Preservation]
\label{lem:positivity}
All tuple entries remain positive under admissible evolution.

Consequently, any increase introduced at one stage propagates
monotonically to all future stages.
\end{lemma}

\begin{proof}
The admissible recursion consists entirely of cumulative additions
of positive quantities.

Terminal duplication copies an existing positive terminal entry,
and no subtraction occurs anywhere in the grammar evolution.

Hence all tuple entries remain positive, and any increase introduced
at one stage propagates to all future descendants.
\end{proof}

% ------------------------------------------------------------

\begin{lemma}[Earlier Duplication Dominance]
\label{lem:earlier-duplication}
Let
\[
J_1,\qquad J_2
\]
be two admissible jump schedules that agree through some index $m$.

Suppose $J_1$ performs a terminal duplication strictly earlier than
$J_2$.

Then the corresponding tuple masses satisfy
\[
|T_n^{(1)}|
\ge
|T_n^{(2)}|
\]
for all sufficiently large $n$, with strict inequality eventually
persistent.
\end{lemma}

\begin{proof}[Heuristic mechanism]
Earlier duplication inserts additional positive mass into the tuple
recursion sooner.

By Lemma~\ref{lem:positivity}, the admissible evolution is monotone:
all future cumulative sums inherit any earlier increase.

Thus earlier duplication permanently amplifies future tuple growth,
yielding eventual domination.
\end{proof}

% ------------------------------------------------------------

\subsection{Canonical Minimal Envelope}

% ------------------------------------------------------------

\begin{theorem}[Canonical Minimal Envelope]
\label{thm:minimal-envelope}
Among all admissible phase continuations of the Sturmian jump grammar,
the canonical intercept-$0$ Beatty phase produces minimal tuple growth.

Equivalently, for every admissible phase shift $\phi$,
\[
|T_n^{(0)}|
\le
|T_n^{(\phi)}|
\]
for all sufficiently large $n$.
\end{theorem}

\begin{proof}[Structural mechanism]
The canonical Beatty phase delays terminal duplication events
as long as admissibly possible.

By Lemma~\ref{lem:earlier-duplication}, any admissible phase shift
that introduces duplication earlier produces permanently larger
future tuple mass.

Hence the canonical phase yields the minimal admissible growth
envelope.
\end{proof}

% ------------------------------------------------------------

\begin{remark}[Minimal Regenerative Orbit]
Exact self-replication of admissible tuples occurs only after
renormalization together with resetting the Beatty phase to the
canonical intercept-$0$ schedule.

However, Theorem~\ref{thm:minimal-envelope} shows that this reset
produces the smallest admissible continuation.

Thus exact regeneration is unnecessary:
the canonical reset already provides a universal lower envelope for
future admissible growth.
\end{remark}

% % ------------------------------------------------------------

% \begin{theorem}[Existence of the Exponential Growth Constant]
% \label{thm:lambda-exists}
% There exists a finite constant
% \[
% \lambda>0
% \]
% such that
% \[
% \lambda
% =
% \lim_{n\to\infty} a(n)^{1/n}.
% \]
% Equivalently,
% \[
% \lim_{n\to\infty}\frac{\log a(n)}{n}
% =
% \log\lambda.
% \]
% \end{theorem}

% \begin{proof}
% Define
% \[
% b(n):=\log a(n).
% \]

% By Theorem~\ref{thm:supermultiplicative},
% \[
% b(n+m)\ge b(n)+b(m),
% \]
% so the sequence $b(n)$ is superadditive.

% Therefore Fekete's lemma applies, yielding existence of the limit
% \[
% \lim_{n\to\infty}\frac{b(n)}{n}
% =
% \sup_n \frac{b(n)}{n}.
% \]

% Exponentiating gives
% \[
% \lim_{n\to\infty} a(n)^{1/n}
% =
% \lambda
% \]
% for some finite constant $\lambda>0$.
% \end{proof}

% % ------------------------------------------------------------

% \begin{remark}[Delayed Compensation and Convexity]
% The canonical Beatty phase minimizes growth by delaying terminal
% duplication events as long as admissibly possible.

% Consequently, local suppression of tuple growth necessarily forces
% larger future compensation once duplication eventually occurs.

% This delayed-compensation mechanism naturally induces:
% \begin{itemize}
% \item block log-convexity,
% \item bounded recovery windows,
% \item asymptotic smoothing of local oscillations.
% \end{itemize}
% \end{remark}

% % ------------------------------------------------------------

% \begin{remark}[Structural Interpretation]
% The admissible division-word grammar defines a positive monotone
% dynamical system controlled by a Sturmian jump schedule.

% The canonical intercept-$0$ Beatty phase acts as the unique minimal
% growth orbit, while all admissible phase shifts dominate it.

% Thus the existence of the exponential scale $\lambda$ emerges from
% an extremal monotone regeneration principle intrinsic to the grammar.
% \end{remark}

% ============================================================
\section{Restricted Duplication Grammar and Spectral Separation}
\label{sec:restricted-grammar}
% ============================================================

The admissible tuple grammar defining \OEIS{A186009} possesses a natural
restricted subsystem obtained by forbidding consecutive terminal
duplication events.

This restricted grammar generates the sequence \OEIS{A047749}, which
therefore serves as a canonical lower-growth sector of the full
admissible grammar.

% ------------------------------------------------------------

\subsection{Restricted Duplication Grammar}

Recall that the admissible grammar evolves through:
\begin{enumerate}
\item cumulative Pascal-type propagation,
\item terminal duplication whenever the jump increment satisfies
\[
\delta(n)=2.
\]
\end{enumerate}

In the full grammar, consecutive duplication events are permitted.

The restricted grammar forbids consecutive duplication, so every
duplication event must be followed by at least one non-duplication step.

Equivalently, if
\[
D
\]
denotes a duplication step and
\[
S
\]
a non-duplication step, then the restricted grammar forbids the local
pattern
\[
DD.
\]

% ------------------------------------------------------------

\begin{definition}[Restricted Duplication Sequence]
Let
\[
a_0(n):=A047749(n).
\]

Then \(a_0(n)\) is generated by the admissible tuple recursion under
the constraint that consecutive terminal duplications are forbidden.
\end{definition}

% ------------------------------------------------------------

\subsection{Initial Coincidence and Spectral Divergence}

The restricted and unrestricted grammars initially coincide because the
canonical Beatty increment sequence contains no consecutive duplication
events at small scales.

Consequently,
\[
A047749(n)=A186009(n+1)
\]
through the first coincidence range.

However, once the unrestricted grammar encounters the first admissible
double--double event, the two systems diverge permanently.

% ------------------------------------------------------------

\begin{lemma}[Permanent Amplification by Consecutive Duplication]
\label{lem:permanent-amplification}
Suppose two admissible tuple evolutions agree through some level \(m\),
and suppose one evolution performs a consecutive terminal duplication
event at level \(m+1\) while the other does not.

Then the unrestricted evolution eventually dominates strictly:
\[
|T_n^{(\mathrm{full})}|
>
|T_n^{(\mathrm{restricted})}|
\]
for all sufficiently large \(n\).
\end{lemma}

\begin{proof}[Structural mechanism]
Terminal duplication inserts additional positive mass into the tuple.

By positivity preservation of the admissible recursion, all future
Pascal-type propagations inherit this excess mass monotonically.

Thus the effect of an earlier duplication event is permanent and
strictly amplifying.
\end{proof}

% \begin{corollary}[Strict Spectral Domination]
% The unrestricted admissible grammar grows asymptotically faster than the
% restricted grammar.

% In particular, if the exponential growth rates exist,
% \[
% \lambda_0
% =
% \limsup_{n\to\infty} A047749(n)^{1/n},
% \]
% and
% \[
% \lambda
% =
% \limsup_{n\to\infty} A186009(n)^{1/n},
% \]
% then
% \[
% \lambda_0 < \lambda.
% \]
% \end{corollary}
\begin{corollary}[Spectral Domination]
If the exponential growth rates exist, then
\[
\lambda_0 \le \lambda.
\]

Moreover, computational evidence and the persistent amplification
mechanism strongly suggest that the inequality is strict:
\[
\lambda_0 < \lambda.
\]
\end{corollary}

% ------------------------------------------------------------
\subsection{Interpretation}

The sequence \OEIS{A047749} therefore represents the minimal
non-consecutive amplification sector of the admissible grammar, while
\OEIS{A186009} represents the full unrestricted positive grammar.

The spectral gap between the two systems is generated entirely by the
presence of consecutive terminal duplication events.

This identifies consecutive duplication as the fundamental local
mechanism responsible for accelerated asymptotic growth of admissible
division-word configurations.

% ------------------------------------------------------------
\section{Grammar-Induced Deficit Compensation}
\label{sec:grammar-compensation}

We formalize a structural mechanism in which jumps in the dyadic
threshold sequence create temporary mass deficits, while the same jumps
simultaneously enlarge the admissible grammar of minimal descent words.
The resulting delayed compensation drives cumulative density toward $1$.

% ------------------------------------------------------------

\begin{definition}[Dyadic Threshold Sequence]
Let
\[
d(n)=A020914(n).
\]
Equivalently,
\[
d(n)=\lfloor n\log_2 3 \rfloor + 1.
\]
This is the minimal dyadic exponent required to exceed the ternary scale
$3^n$.
\end{definition}

% ------------------------------------------------------------

\begin{definition}[Jump Index]
An index $n\ge0$ is called a \emph{jump index} if
\[
d(n+1)-d(n)=2.
\]
Otherwise,
\[
d(n+1)-d(n)=1.
\]
\end{definition}

% ------------------------------------------------------------

\begin{lemma}[Binary Step Law]
\label{lem:binary-step}
For every $n\ge0$,
\[
d(n+1)-d(n)\in\{1,2\}.
\]
Hence the sequence $d(n)$ evolves only by unit steps or jump steps.
\end{lemma}

\begin{proof}
Using
\[
d(n)=\lfloor n\log_2 3\rfloor+1,
\]
we obtain
\[
d(n+1)-d(n)
=
\lfloor (n+1)\log_2 3\rfloor-\lfloor n\log_2 3\rfloor.
\]
Since
\[
1<\log_2 3<2,
\]
the floor can increase only by $1$ or $2$.
\end{proof}

% ------------------------------------------------------------

\begin{lemma}[Single Dyadic Deficit Law]
\label{lem:deficit-law}
If $n$ is a jump index, then exactly one intermediate dyadic level is
skipped relative to unit growth.

Equivalently, the denominator factor changes from the expected scale
$2^{d(n)+1}$ to the actual scale $2^{d(n)+2}$, reducing the associated
mass contribution by a factor of $1/2$.
\end{lemma}

\begin{proof}
At a jump index,
\[
d(n+1)=d(n)+2.
\]
Relative to the unit-step scale $d(n)+1$, one additional factor of $2$
appears in the denominator, so the contribution is halved.
\end{proof}

% ------------------------------------------------------------

\begin{definition}[Minimal Descent Mass]
Let
\[
N(n)=\frac{A186009(n+1)}{2^{d(n)}}.
\]
This is the total density contribution of new minimal descent words at
level $n$.
\end{definition}

% ------------------------------------------------------------

\begin{lemma}[Jump Enlarges the Admissible Grammar]
\label{lem:jump-grammar}
At every jump index, the admissible prefix-sum region defining minimal
descent words strictly enlarges at subsequent levels.

Consequently, new minimal descent words must occur after each jump.
\end{lemma}

\begin{proof}[Proof Sketch]
The admissible words are determined by prefix constraints of the form
\[
H(i)<d(i+1), \qquad H(m-1)=d(m).
\]
When $d(n)$ jumps, the terminal admissible sum increases while prior
prefix constraints remain unchanged. This creates new feasible lattice
points in the constrained composition region, yielding new admissible
words.
\end{proof}

% ------------------------------------------------------------

\begin{lemma}[Lagged Numerator Response]
\label{lem:lagged-response}
There exists an integer $L\ge1$ such that every jump at index $n$
produces at least one new minimal descent word among the levels
\[
n+1,n+2,\dots,n+L.
\]
\end{lemma}

\begin{proof}[Heuristic Basis]
Computational enumeration shows that new admissible words appear after
every jump, with uniformly bounded delay. The bound $L$ is empirically
small.
\end{proof}

% ------------------------------------------------------------

\begin{lemma}[Finite Compensation Packet]
\label{lem:packet}
There exist constants $L\ge1$ and $c>0$ such that every jump index $n$
satisfies
\[
\sum_{i=1}^{L} N(n+i)\ge c\,2^{-d(n)}.
\]
That is, each jump-induced unit deficit is compensated by a bounded
future packet of new descent mass.
\end{lemma}

% ------------------------------------------------------------

\begin{definition}[Residual Gap]
Let
\[
C(m)=\sum_{k=0}^{m}N(k),
\qquad
R(m)=1-C(m).
\]
\end{definition}

% ------------------------------------------------------------

\begin{theorem}[Compensation Implies Gap Contraction]
\label{thm:gap-contraction}
Suppose the finite compensation packet property holds uniformly.
Then there exist constants $M\ge1$ and $\epsilon>0$ such that
\[
R(m+M)\le (1-\epsilon)R(m)
\]
for all sufficiently large $m$.
\end{theorem}

\begin{proof}[Proof Sketch]
Each jump creates at most one unit deficit by
Lemma~\ref{lem:deficit-law}. By Lemma~\ref{lem:packet}, each such deficit
is repaired within bounded delay by a positive packet of future mass.
Uniform boundedness implies a fixed positive fraction of the remaining
gap is removed over each window of length $M$.
\end{proof}

% ------------------------------------------------------------

\begin{corollary}[Density Exhaustion]
\label{cor:density-exhaustion}
Under the hypotheses of Theorem~\ref{thm:gap-contraction},
\[
\lim_{m\to\infty} C(m)=1.
\]
\end{corollary}

\begin{proof}
Iterating the contraction inequality gives
\[
R(m+kM)\le (1-\epsilon)^kR(m)\to0.
\]
Hence $C(m)\to1$.
\end{proof}

% ------------------------------------------------------------

\begin{remark}[Critical Reduction]
The remaining task is purely combinatorial: prove a uniform bound on the
delay and size of compensation packets generated by jump indices in the
admissible division-word grammar.
\end{remark}

% ============================================================
\clearpage
\bibliographystyle{unsrtnat}
\bibliography{src/latex/references}

\end{document}